
\section{Related Work}

\subsection{Low-Rank Matrix Decomposition}

\subsubsection{SVD and Eckart--Young}

The truncated SVD is optimal in Frobenius norm (Eckart \& Young, 1936). However, 
the factors $U, V$ are dense and orthogonal, which limits interpretability.

\subsubsection{CUR Decomposition}

Mahoney \& Drineas (2009) introduced CUR as an interpretable alternative: 
select $r$ columns and $r$ rows of $W$, then reconstruct as $C \cdot U^+ \cdot R$ 
where $U$ is the intersection. The key insight is that **statistical leverage 
scores** (Mahoney et al., 2012) provide the optimal sampling probabilities. 
PHANTOM uses deterministic top-$r$ selection by leverage score.

\subsubsection{Randomized SVD}

Halko et al. (2011) provide randomized SVD with complexity 
$\mathcal{O}(d^2 r)$. We use exact SVD for leverage computation because 
matrix sizes are manageable.

\subsection{Sparse Attention}

\subsubsection{Fixed-Pattern Sparsity}

Longformer (Beltagy et al., 2020) uses a sliding window plus global tokens. 
BigBird (Zaheer et al., 2020) adds random edges. Both use \textbf{fixed} patterns 
that do not adapt to input.

\subsubsection{Content-Based Sparsity}

Reformer (Kitaev et al., 2020) uses LSH to group similar queries and keys. 
Routing Transformer (Roy et al., 2021) uses online k-means. PHANTOM uses 
\textbf{multi-round angular LSH} with a rectangular hash function that is 
differentiable through sign operations.

\subsubsection{Comparisons}

\begin{table}[H]
\centering
\small
\begin{tabular}{llll}
\toprule
\textbf{Method} & \textbf{Pattern} & \textbf{Adaptive} & \textbf{Complexity} \\
\midrule
Longformer & Sliding window + global & No & $\mathcal{O}(L \sqrt{L})$ \\
BigBird & Window + random + global & No & $\mathcal{O}(L)$ \\
Reformer & LSH buckets & Yes & $\mathcal{O}(L \log L)$ \\
Routing & k-means clusters & Yes & $\mathcal{O}(L \log L)$ \\
\textbf{PHANTOM} & Multi-round angular LSH & Yes & $\mathcal{O}(L \cdot B)$ \\
\bottomrule
\end{tabular}
\end{table}

\subsection{Attention Activations}

\subsubsection{Softmax}

Standard softmax has two limitations: (i) it never reaches exactly zero, so 
far-tail positions receive non-negligible weight; (ii) it is sensitive to 
outliers.

\subsubsection{Sparsemax}

Martins \& Astudillo (2016) propose sparsemax, which projects onto the simplex 
and produces sparse attention. However, sparsemax has no closed-form for 
batched computation.

\subsubsection{Softpick}

Lu et al. (2024) propose softpick: $\text{ReLU}(e^{(x_i - m)} - e^{-m})$ 
where $m = \max_j x_j$. This is sparse and has closed form. We use a 
\textbf{rectified} variant with a softmax fallback when the denominator is zero.

\subsection{Post-Training Compression}

\begin{table}[H]
\centering
\small
\begin{tabular}{lllll}
\toprule
\textbf{Method} & \textbf{Type} & \textbf{Retrain?} & \textbf{Compression} & \textbf{Attention?} \\
\midrule
SVD (Eckart-Young) & Low-rank & No & Ratio $r$ & No \\
A3 & Fisher-weighted SVD & LoRA & Ratio $r$ & No \\
AASVD & Covariance-weighted SVD & LoRA & Ratio $r$ & No \\
SparseGPT & Pruning & No & 50\% & No \\
\textbf{PHANTOM} & CUR + LSH & LoRA & Ratio $r$ & \textbf{Yes} \\
\bottomrule
\end{tabular}
\end{table}

\textbf{Key distinction.} PHANTOM is the only method in this list that compresses 
both \textbf{projections} (via CUR) \textbf{and attention} (via LSH).

\subsection{Position of Our Work}

PHANTOM is unique in combining:
\begin{enumerate}
\item \textbf{CUR decomposition} for interpretable low-rank projections.
\item \textbf{Multi-round LSH} for content-adaptive sparse attention.
\item \textbf{Rectified softpick} for stable sparse attention.

\end{enumerate}
Prior work addresses each of these individually, but no prior method combines 
all three.
\begin{lstlisting}

---
\end{lstlisting}